\documentclass[11pt]{article}
\usepackage[margin=1in]{geometry}
\usepackage{times,amsmath,amssymb,booktabs,graphicx,adjustbox,url,natbib}
\usepackage{hyperref}
\emergencystretch=3em
\title{Scaling and reading interpretable spline causal discovery:\\
adaptive sparsity, parent screening, and shape-faithful edge functions}
\author{[authors to be confirmed]}
\date{}
\begin{document}
\maketitle

\begin{abstract}
SPADE learns a non-linear additive-noise DAG whose edges are B-spline functions, so every learned
dependency can be plotted and inspected. It has three practical weaknesses: its group-lasso weight must
shrink with the number of variables $d$, it shrinks strong edges as much as weak ones, and its raw spline
curves wiggle, which hides the shapes it is meant to expose. We address each. (i) An adaptive group-lasso
averaged over half-sample fits (SPADE-AB) raises instantaneous-DAG AUROC by $0.014$ on fresh seeds
(Wilcoxon $p=0.004$, $n=45$ runs). (ii) A width-aware penalty $\lambda_g=\min(0.01,0.2/d)$ with
per-node parent screening keeps AUROC at $0.96$ to $0.98$ for $d=50$ to $200$, where the fixed
penalty collapses to $0.58$ to $0.82$. (iii) A smoothness-penalized refit of the retained edges lifts the
agreement between learned and true curve class from $0.12$--$0.25$ to $0.93$--$0.94$ without changing the
recovered structure. On the Sachs protein-signalling network, the one real benchmark with a ground-truth graph that
we use, SPADE-AB attains AUROC $0.664$ against at most $0.566$ for the other methods tested, although absolute
accuracy remains low. We state the limits: all comparisons at $d\ge100$ are against linear baselines
(non-linear competitors do not run at that size), and the synthetic gains are modest.
\end{abstract}

\section{Introduction}
SPADE \citep{Dang2026spade} couples Kolmogorov--Arnold-style B-spline edge functions \citep{Liu2024} with
an acyclicity constraint \citep{Zheng2018} and a group-lasso, giving a non-linear causal graph whose edges
are readable curves. Its published evaluation stopped at $d=50$ and used one fixed sparsity weight. This
paper asks three questions a practitioner would ask next. Does it work at hundreds of variables? Can the
sparsity step be made less biased? And are the curves it draws actually faithful to the mechanisms that
generated the data, which is the property that justifies calling the model interpretable? We answer with
protocols fixed before the final runs, candidates chosen on validation data disjoint from the evaluation
data, and failures reported.

\section{Method}
\paragraph{Base model.} Each variable $x_i$ is predicted only through spline edges of its candidate
causes, $\hat x_i=b_i+\sum_{j\ne i}\phi_{ij}(x_j)$, trained with MSE, a group-lasso
$\lambda_g\sum_{ij}\lVert\theta_{ij}\rVert_2$ over the coefficient vector of each edge, and the NOTEARS
penalty on the edge-norm matrix \citep{Zheng2018}. Edge scores are the standard deviation of each edge's
contribution.

\paragraph{SPADE-AB (adaptive sparsity with averaging).} Stage 1 is the base fit. Stage 2 refits with
weights $w_{ij}=(\lVert\theta_{ij}\rVert+\epsilon)^{-1}$ normalized to unit mean (adaptive group-lasso,
\citealt{Zou2006}), which penalizes weak edges more and strong edges less. The final score is the mean of
max-normalized scores over five fits on random half samples. On validation seeds neither the adaptive step
nor the averaging alone was reliably better; the combination had the best mean AUROC at all three widths.
Plain adaptive fitting was no better than the base fit and refitting without penalty was worse
(Table~\ref{tab:valsel}).

\paragraph{Width-aware penalty and parent screening.} The group-lasso sums over $d(d-1)$ edges, so a fixed
$\lambda_g$ over-penalizes as $d$ grows. We use $\lambda_g(d)=\min(0.01,c/d)$, and before fitting keep for
each node only the $K$ candidate causes with the highest ExtraTrees importance, masking all other edges.
$c=0.2$ and $K=15$ were chosen on validation graphs at $d=50$ and $100$.

\paragraph{Shape-faithful display refit.} Raw spline coefficients are unconstrained at the grid tails and
produce spurious wiggles. After structure is fixed, we refit the retained edges (top $2d$ by score) with a
second-difference penalty $0.01\sum\lVert D^2\theta_{ij}\rVert^2$ and plot those curves. Structure
(ranking) always comes from the unsmoothed fit, so the display step cannot change which edges are found.

\section{Experiments}
Instantaneous data come from a non-linear additive-noise generator with non-Gaussian noise (mechanisms
$\tanh$, $\sin$, quadratic, rational). Validation and evaluation seeds are disjoint throughout. Baselines are
run at their reference settings as in the original SPADE study.

\subsection{Sparsity: SPADE-AB}
\begin{table}[t]\centering\small
\caption{Fresh seeds (300--319 for $d=6,10$; 300--304 for $d=20$): instantaneous-DAG recovery.}
\label{tab:ab}
\begin{tabular}{lcccccc}
\toprule
 & \multicolumn{2}{c}{$d{=}6$} & \multicolumn{2}{c}{$d{=}10$} & \multicolumn{2}{c}{$d{=}20$} \\
 Method & AUROC & AUPRC & AUROC & AUPRC & AUROC & AUPRC \\ \midrule
SPADE & 0.961 & 0.924 & 0.948 & 0.867 & 0.914 & 0.800 \\
SPADE-AB & 0.974 & 0.966 & 0.961 & 0.922 & 0.935 & 0.848 \\
\bottomrule
\end{tabular}

\end{table}
Table~\ref{tab:ab} shows consistent gains. Pooled over the 45 runs the paired difference is
+0.014 AUROC, with 24 wins, 13 ties and 8 losses over 45 runs (Wilcoxon $p=0.0039$). The gain per width is small and
not individually significant; only the pooled test is. In an earlier comparison SPADE trailed NoGAM at
$d=10$ (0.926 vs 0.955); with SPADE-AB the gap on the same seeds narrows from 0.063 to 0.015 but does not close.
\emph{Environment note.} Re-running the base SPADE in a newer library environment gave $d=10$ AUROC $0.892$ on the
official seeds, below the originally reported $0.926$. Single-run variance across library versions is therefore
about $\pm0.03$ at $d=10$, comparable to some of the gains reported here, which is why we rely on paired,
same-environment comparisons.

\begin{table}[t]\centering\small
\caption{Validation selection (seeds 200--219, $d=6,10$; 200--204, $d=20$): mean AUROC.}
\label{tab:valsel}
\begin{tabular}{lccc}\toprule Variant & $d{=}6$ & $d{=}10$ & $d{=}20$\\\midrule
base & 0.958 & 0.923 & 0.871\\ adaptive & 0.952 & 0.917 & 0.871\\ refit (no penalty) & 0.930 & 0.900 & 0.855\\
bootstrap average & 0.971 & 0.939 & 0.877\\ adaptive + average (SPADE-AB) & 0.974 & 0.947 & 0.893\\\bottomrule\end{tabular}
\end{table}

\subsection{Scale: $d=50$ to $200$}
\begin{table}[t]\centering\small
\caption{Fresh seeds 600--602, $n=3{,}000$. SPADE-scaled: $\lambda_g=\min(0.01,0.2/d)$ and parent screening
($K=15$). GOLEM and NOTEARS-MLP were run only at $d=50$ (each fit exceeded 15 minutes); non-linear competitors
(NoGAM, DAGMA-nonlinear, GraN-DAG) cannot run at $d\ge100$ and SCORE did not finish within our budget.}
\label{tab:scale}
\adjustbox{max width=\textwidth}{\begin{tabular}{llccc}
\toprule
$d$ & Method & AUROC & AUPRC & seeds \\ \midrule
50 & DAGMA-linear & 0.648$\pm$0.046 & 0.138 & 3 \\
50 & GOLEM & 0.807$\pm$0.043 & 0.188 & 3 \\
50 & NOTEARS-MLP & 0.530$\pm$0.050 & 0.023 & 3 \\
50 & NOTEARS-linear & 0.801$\pm$0.018 & 0.355 & 3 \\
50 & SPADE ($\lambda_g{=}0.01$) & 0.815$\pm$0.039 & 0.438 & 3 \\
50 & SPADE-scaled & 0.961$\pm$0.011 & 0.774 & 3 \\
100 & DAGMA-linear & 0.655$\pm$0.031 & 0.116 & 3 \\
100 & NOTEARS-linear & 0.578$\pm$0.031 & 0.069 & 3 \\
100 & SPADE ($\lambda_g{=}0.01$) & 0.678$\pm$0.026 & 0.043 & 3 \\
100 & SPADE-scaled & 0.982$\pm$0.009 & 0.830 & 3 \\
200 & DAGMA-linear & 0.674$\pm$nan & 0.155 & 1 \\
200 & NOTEARS-linear & 0.471$\pm$nan & 0.007 & 1 \\
200 & SPADE ($\lambda_g{=}0.01$) & 0.579$\pm$nan & 0.006 & 1 \\
200 & SPADE-scaled & 0.979$\pm$nan & 0.766 & 1 \\
\bottomrule
\end{tabular}
}
\end{table}
Table~\ref{tab:scale} shows the fixed penalty collapsing as $d$ grows, while the width-aware rule with screening
holds AUROC near $0.96$--$0.98$. The comparison at $d\ge100$ is only against linear methods, which are
mis-specified for this data; it shows SPADE-scaled works at this size, not that it beats the best non-linear learner.
Run times are not reported because jobs shared a four-core machine.

\subsection{Interpretability, measured}
For every true edge that a method ranks within its top-$|E|$ scored edges, we compare the learned curve with the
true mechanism: $R^2$ of the best affine map, and agreement of the qualitative class (increasing, decreasing,
non-monotone). ``Strongly non-linear'' edges are those a straight line explains less than half of (sin and quadratic).
\begin{table}[t]\centering\small
\caption{Shape fidelity on fresh seeds 400--419. MLP-edge: same bottleneck, penalty and budget with $1\!\to\!3\!\to\!1$
$\tanh$ edges. SPADE-2stage: structure from the plain SPADE fit, curves from the smoothed refit.}
\label{tab:fid}
\adjustbox{max width=\textwidth}{\begin{tabular}{llcccc}
\toprule
 & & \multicolumn{2}{c}{found edges (all)} & \multicolumn{2}{c}{found, strongly non-linear} \\
$d$ & Method & $R^2$ & class agr. & $R^2$ & ($n$) \\ \midrule
6 & MLP-edge & 0.92 & 0.91 & 0.23 & (4) \\
6 & SPADE & 0.91 & 0.25 & 0.71 & (16) \\
6 & SPADE-2stage & 0.96 & 0.93 & 0.85 & (16) \\
10 & MLP-edge & 0.95 & 0.94 & 0.24 & (3) \\
10 & SPADE & 0.91 & 0.12 & 0.55 & (15) \\
10 & SPADE-2stage & 0.96 & 0.94 & 0.63 & (15) \\
\bottomrule
\end{tabular}
}
\end{table}
Table~\ref{tab:fid} has two messages. First, raw SPADE curves recover the right values ($R^2\approx0.91$) but the wrong
qualitative class on most edges, because of wiggles; the smoothed display refit repairs this ($0.93$--$0.94$) and
reaches the best $R^2$ ($0.96$). Second, on strongly non-linear edges, SPADE finds many more of them than the
matched-capacity MLP edges and recovers their shape far better. The ``found'' rule uses the true number of
edges, so it favours no method but is an oracle count. A metric computed in an early run was wrong because of an
indexing error (the curve was read from the wrong input column); we fixed it and re-ran before any result was used.

\subsection{A real benchmark with ground truth: Sachs}
The Sachs data are 7{,}466 single-cell measurements of 11 phospho-proteins with an 18-edge consensus graph. We use
10 random subsamples of 2{,}000 cells after a log transform; no method is tuned on these data.
\begin{table}[t]\centering\small
\caption{Sachs, 10 subsamples of $2{,}000$ cells (chance AUPRC $=0.164$). SPADE-AB and plain SPADE use the settings
fixed on synthetic data.}
\label{tab:sachs}
\begin{tabular}{lcccc}
\toprule
Method & AUROC & AUPRC & top-18 hits & SHD \\ \midrule
SPADE-AB & 0.664$\pm$0.043 & 0.319 & 6.0 & 24.0 \\
SPADE & 0.601$\pm$0.048 & 0.289 & 5.8 & 24.4 \\
NOTEARS-linear & 0.566$\pm$0.015 & 0.239 & 5.0 & 26.0 \\
GOLEM & 0.528$\pm$0.037 & 0.195 & 3.6 & 28.8 \\
DAGMA-nonlinear & 0.486$\pm$0.117 & 0.238 & 3.8 & 28.4 \\
GraN-DAG & 0.484$\pm$0.002 & 0.164 & 2.0 & 32.0 \\
NoGAM & 0.465$\pm$0.040 & 0.157 & 6.2 & 50.0 \\
SCORE & 0.461$\pm$0.037 & 0.156 & 6.0 & 49.8 \\
DAGMA-linear & 0.432$\pm$0.035 & 0.167 & 3.0 & 30.0 \\
NOTEARS-MLP & 0.297$\pm$0.036 & 0.117 & 0.1 & 35.8 \\
\bottomrule
\end{tabular}

\end{table}
SPADE-AB is best (Table~\ref{tab:sachs}); it beats every other method in 8 or more of 10 resamples ($p\le0.010$).
Absolute accuracy is low: it recovers about 6 of the 18 edges among its top 18. Sachs is difficult for all methods,
and this result shows relative rather than good performance.

\section{Limitations}
(1) Gains from SPADE-AB are modest. (2) At $d\ge100$ only linear baselines were run. (3) The parent screening step
can discard true parents, so recall is bounded by the screen. (4) The interpretability metric uses synthetic
mechanisms and an oracle edge count; it is not a human study. (5) Sachs is a single real data set. (6) On the
financial and weather series we tried, no learned one-step forecaster beat a naive last-value prediction, SPADE
variants included, so we make no forecasting claim.

\section{Conclusion}
A width-aware penalty with parent screening makes SPADE usable at $d=200$; adaptive, averaged sparsity gives a small
but consistent accuracy gain; and a smoothed display refit makes the learned curves match the true mechanism's class
without altering the recovered graph. The next step is a comparison with non-linear learners at $d\ge100$ and a
human study of whether the curves change decisions.

\bibliographystyle{plainnat}
\bibliography{references}
\end{document}
